% Options for packages loaded elsewhere
\PassOptionsToPackage{unicode}{hyperref}
\PassOptionsToPackage{hyphens}{url}
\PassOptionsToPackage{dvipsnames,svgnames,x11names}{xcolor}
%
\documentclass[
  letterpaper,
  DIV=11,
  numbers=noendperiod]{scrartcl}

\usepackage{amsmath,amssymb}
\usepackage{iftex}
\ifPDFTeX
  \usepackage[T1]{fontenc}
  \usepackage[utf8]{inputenc}
  \usepackage{textcomp} % provide euro and other symbols
\else % if luatex or xetex
  \usepackage{unicode-math}
  \defaultfontfeatures{Scale=MatchLowercase}
  \defaultfontfeatures[\rmfamily]{Ligatures=TeX,Scale=1}
\fi
\usepackage{lmodern}
\ifPDFTeX\else  
    % xetex/luatex font selection
\fi
% Use upquote if available, for straight quotes in verbatim environments
\IfFileExists{upquote.sty}{\usepackage{upquote}}{}
\IfFileExists{microtype.sty}{% use microtype if available
  \usepackage[]{microtype}
  \UseMicrotypeSet[protrusion]{basicmath} % disable protrusion for tt fonts
}{}
\makeatletter
\@ifundefined{KOMAClassName}{% if non-KOMA class
  \IfFileExists{parskip.sty}{%
    \usepackage{parskip}
  }{% else
    \setlength{\parindent}{0pt}
    \setlength{\parskip}{6pt plus 2pt minus 1pt}}
}{% if KOMA class
  \KOMAoptions{parskip=half}}
\makeatother
\usepackage{xcolor}
\setlength{\emergencystretch}{3em} % prevent overfull lines
\setcounter{secnumdepth}{-\maxdimen} % remove section numbering
% Make \paragraph and \subparagraph free-standing
\makeatletter
\ifx\paragraph\undefined\else
  \let\oldparagraph\paragraph
  \renewcommand{\paragraph}{
    \@ifstar
      \xxxParagraphStar
      \xxxParagraphNoStar
  }
  \newcommand{\xxxParagraphStar}[1]{\oldparagraph*{#1}\mbox{}}
  \newcommand{\xxxParagraphNoStar}[1]{\oldparagraph{#1}\mbox{}}
\fi
\ifx\subparagraph\undefined\else
  \let\oldsubparagraph\subparagraph
  \renewcommand{\subparagraph}{
    \@ifstar
      \xxxSubParagraphStar
      \xxxSubParagraphNoStar
  }
  \newcommand{\xxxSubParagraphStar}[1]{\oldsubparagraph*{#1}\mbox{}}
  \newcommand{\xxxSubParagraphNoStar}[1]{\oldsubparagraph{#1}\mbox{}}
\fi
\makeatother

\usepackage{color}
\usepackage{fancyvrb}
\newcommand{\VerbBar}{|}
\newcommand{\VERB}{\Verb[commandchars=\\\{\}]}
\DefineVerbatimEnvironment{Highlighting}{Verbatim}{commandchars=\\\{\}}
% Add ',fontsize=\small' for more characters per line
\usepackage{framed}
\definecolor{shadecolor}{RGB}{241,243,245}
\newenvironment{Shaded}{\begin{snugshade}}{\end{snugshade}}
\newcommand{\AlertTok}[1]{\textcolor[rgb]{0.68,0.00,0.00}{#1}}
\newcommand{\AnnotationTok}[1]{\textcolor[rgb]{0.37,0.37,0.37}{#1}}
\newcommand{\AttributeTok}[1]{\textcolor[rgb]{0.40,0.45,0.13}{#1}}
\newcommand{\BaseNTok}[1]{\textcolor[rgb]{0.68,0.00,0.00}{#1}}
\newcommand{\BuiltInTok}[1]{\textcolor[rgb]{0.00,0.23,0.31}{#1}}
\newcommand{\CharTok}[1]{\textcolor[rgb]{0.13,0.47,0.30}{#1}}
\newcommand{\CommentTok}[1]{\textcolor[rgb]{0.37,0.37,0.37}{#1}}
\newcommand{\CommentVarTok}[1]{\textcolor[rgb]{0.37,0.37,0.37}{\textit{#1}}}
\newcommand{\ConstantTok}[1]{\textcolor[rgb]{0.56,0.35,0.01}{#1}}
\newcommand{\ControlFlowTok}[1]{\textcolor[rgb]{0.00,0.23,0.31}{\textbf{#1}}}
\newcommand{\DataTypeTok}[1]{\textcolor[rgb]{0.68,0.00,0.00}{#1}}
\newcommand{\DecValTok}[1]{\textcolor[rgb]{0.68,0.00,0.00}{#1}}
\newcommand{\DocumentationTok}[1]{\textcolor[rgb]{0.37,0.37,0.37}{\textit{#1}}}
\newcommand{\ErrorTok}[1]{\textcolor[rgb]{0.68,0.00,0.00}{#1}}
\newcommand{\ExtensionTok}[1]{\textcolor[rgb]{0.00,0.23,0.31}{#1}}
\newcommand{\FloatTok}[1]{\textcolor[rgb]{0.68,0.00,0.00}{#1}}
\newcommand{\FunctionTok}[1]{\textcolor[rgb]{0.28,0.35,0.67}{#1}}
\newcommand{\ImportTok}[1]{\textcolor[rgb]{0.00,0.46,0.62}{#1}}
\newcommand{\InformationTok}[1]{\textcolor[rgb]{0.37,0.37,0.37}{#1}}
\newcommand{\KeywordTok}[1]{\textcolor[rgb]{0.00,0.23,0.31}{\textbf{#1}}}
\newcommand{\NormalTok}[1]{\textcolor[rgb]{0.00,0.23,0.31}{#1}}
\newcommand{\OperatorTok}[1]{\textcolor[rgb]{0.37,0.37,0.37}{#1}}
\newcommand{\OtherTok}[1]{\textcolor[rgb]{0.00,0.23,0.31}{#1}}
\newcommand{\PreprocessorTok}[1]{\textcolor[rgb]{0.68,0.00,0.00}{#1}}
\newcommand{\RegionMarkerTok}[1]{\textcolor[rgb]{0.00,0.23,0.31}{#1}}
\newcommand{\SpecialCharTok}[1]{\textcolor[rgb]{0.37,0.37,0.37}{#1}}
\newcommand{\SpecialStringTok}[1]{\textcolor[rgb]{0.13,0.47,0.30}{#1}}
\newcommand{\StringTok}[1]{\textcolor[rgb]{0.13,0.47,0.30}{#1}}
\newcommand{\VariableTok}[1]{\textcolor[rgb]{0.07,0.07,0.07}{#1}}
\newcommand{\VerbatimStringTok}[1]{\textcolor[rgb]{0.13,0.47,0.30}{#1}}
\newcommand{\WarningTok}[1]{\textcolor[rgb]{0.37,0.37,0.37}{\textit{#1}}}

\providecommand{\tightlist}{%
  \setlength{\itemsep}{0pt}\setlength{\parskip}{0pt}}\usepackage{longtable,booktabs,array}
\usepackage{calc} % for calculating minipage widths
% Correct order of tables after \paragraph or \subparagraph
\usepackage{etoolbox}
\makeatletter
\patchcmd\longtable{\par}{\if@noskipsec\mbox{}\fi\par}{}{}
\makeatother
% Allow footnotes in longtable head/foot
\IfFileExists{footnotehyper.sty}{\usepackage{footnotehyper}}{\usepackage{footnote}}
\makesavenoteenv{longtable}
\usepackage{graphicx}
\makeatletter
\def\maxwidth{\ifdim\Gin@nat@width>\linewidth\linewidth\else\Gin@nat@width\fi}
\def\maxheight{\ifdim\Gin@nat@height>\textheight\textheight\else\Gin@nat@height\fi}
\makeatother
% Scale images if necessary, so that they will not overflow the page
% margins by default, and it is still possible to overwrite the defaults
% using explicit options in \includegraphics[width, height, ...]{}
\setkeys{Gin}{width=\maxwidth,height=\maxheight,keepaspectratio}
% Set default figure placement to htbp
\makeatletter
\def\fps@figure{htbp}
\makeatother

\KOMAoption{captions}{tableheading}
\makeatletter
\@ifpackageloaded{caption}{}{\usepackage{caption}}
\AtBeginDocument{%
\ifdefined\contentsname
  \renewcommand*\contentsname{Table of contents}
\else
  \newcommand\contentsname{Table of contents}
\fi
\ifdefined\listfigurename
  \renewcommand*\listfigurename{List of Figures}
\else
  \newcommand\listfigurename{List of Figures}
\fi
\ifdefined\listtablename
  \renewcommand*\listtablename{List of Tables}
\else
  \newcommand\listtablename{List of Tables}
\fi
\ifdefined\figurename
  \renewcommand*\figurename{Figure}
\else
  \newcommand\figurename{Figure}
\fi
\ifdefined\tablename
  \renewcommand*\tablename{Table}
\else
  \newcommand\tablename{Table}
\fi
}
\@ifpackageloaded{float}{}{\usepackage{float}}
\floatstyle{ruled}
\@ifundefined{c@chapter}{\newfloat{codelisting}{h}{lop}}{\newfloat{codelisting}{h}{lop}[chapter]}
\floatname{codelisting}{Listing}
\newcommand*\listoflistings{\listof{codelisting}{List of Listings}}
\makeatother
\makeatletter
\makeatother
\makeatletter
\@ifpackageloaded{caption}{}{\usepackage{caption}}
\@ifpackageloaded{subcaption}{}{\usepackage{subcaption}}
\makeatother

\ifLuaTeX
  \usepackage{selnolig}  % disable illegal ligatures
\fi
\usepackage{bookmark}

\IfFileExists{xurl.sty}{\usepackage{xurl}}{} % add URL line breaks if available
\urlstyle{same} % disable monospaced font for URLs
\hypersetup{
  pdftitle={Basic Description of Models},
  colorlinks=true,
  linkcolor={blue},
  filecolor={Maroon},
  citecolor={Blue},
  urlcolor={Blue},
  pdfcreator={LaTeX via pandoc}}


\title{Basic Description of Models}
\author{}
\date{}

\begin{document}
\maketitle


\section{Background}\label{background}

Data are collected from FlowCam images of microplankton from the
Mission-Aransas Estuary. Samples range from 2014 to 2021 collected
roughly every month at four stations. Images were initially segemented
using FlowCam's VisualSpreadsheet (version ranged from 2.0 to 4.0) which
provided collages of individual vignettes of particles. Single image
files were then extracted using MorphoCut {[}will need to cite and link
code{]}, which formats images to be individually processed in EcoTaxa.
The vignettes of particles were sorted manually to create a learning set
{[}note the total size later{]} which then was used to train and predict
classifications for the other vignettes. Ultimately, all particles with
ABD {[}check for later{]} greater than 20 microns were manually
validated by a single expert and that is the data used here. {[}will
need to reference ID guides{]}

In imaging, VisualSpreadsheet reports several metrics related to size
including length and width, which were used to calculate biovolume of
each individual. The biovolume of each individual can then be multiplied
to get a biomass (in carbon). The total biomass of any group of plankton
per unit volume can then be calculated:

\begin{align}
mgCm^{-3} = \frac{mgC}{v}
\end{align}

In plankton ecology/biological oceanography this is what researchers are
typically trying to report. Estimates of \(mgCm^{-3}\) are either done
with microscopy or imaging but in either case, it involves converting
counts to biovolume to biomass. Because the counts of individuals tend
to have large amounts of zeros.

There are several ways to approach this:

\section{Approach 1: Log-normal
regression}\label{approach-1-log-normal-regression}

An initial approach may be to simply model the data using a log-normal
distribution. This is typically the primary approach taken:

\subsubsection{Basic Case:}\label{basic-case}

Terms:

\begin{itemize}
\item
  \(y\) - biomass concentration (calculated from data of biovolume /
  measured volume)
\item
  \(i\) - the number of observations
\item
  \(\sigma\) - sampling error / noise
\end{itemize}

\begin{align}
    \text{Data Model:}\\
    log(y_i) &\sim N(\mu_i, \sigma^2)\\
    \\
    \text{Parameter Model:}\\
    \mu_i &= X'\beta
\end{align}

\begin{align}
    y_{g,i} &\sim LN(\mu_{g,i}, s^2)\\
    \\
    \mu_{g,i} &= \pmb{\beta_g X_i}\\
    \\
    \pmb{\beta_g} &\sim N(\pmb{\theta_g}, \sigma^2_g)\\
    \pmb{\theta_g} &\sim N(0, \pmb{\tau^2})\\
    \pmb{\tau} &\sim [\pmb{\tau}]
\end{align}

\$\$

\section{Approach 1: Hurdle Model}\label{approach-1-hurdle-model}

The 0's are an issue given that there are many instances of
non-observation where biomass is zero (e.g.~\(y_{g,i} = 0\)) For this a
hurdle model may be appropriate as we can separately model the occurance
probability and Such as:

\[
\begin{align}
    y_{i} \sim \ &\left\{ 
        \begin{array}{l}
            0; z_i = 0\\
            d_i; z_i = 1
        \end{array} \right.
\end{align}
\]

Now there's two models which contribute to the observed biomass \(y_i\).
First, we can model whether or not a given plankter is present based on
\(z_i\)

\begin{Shaded}
\begin{Highlighting}[]

\NormalTok{\textbackslash{}begin\{align\}}
\NormalTok{    \&\textbackslash{}text\{Occurance Model:\}\textbackslash{}\textbackslash{}}
\NormalTok{    z\_i \&\textbackslash{}sim Bern(\textbackslash{}pi\_i) \textbackslash{}\textbackslash{} }
\NormalTok{    logit(\textbackslash{}pmb\{\textbackslash{}pi\}) \&= \textbackslash{}pmb\{X\}\textquotesingle{}\textbackslash{}pmb\{\textbackslash{}alpha\}\textbackslash{}\textbackslash{}}
    
\NormalTok{    \textbackslash{}pmb\{\textbackslash{}alpha\} \&\textbackslash{}sim [\textbackslash{}pmb\{\textbackslash{}alpha\}]}
\NormalTok{\textbackslash{}end\{align\}}
\end{Highlighting}
\end{Shaded}

Here, \(\pmb{\alpha}\) denotes the coefficient matrix for the effect of
different environmental covariates.

Then separately, we can model the expected biomass concentration, given
that some plankton are observed (\(z_i = 1\))

\$\$ \begin{align}
    &\text{Biomass Model:}\\
    d_i &\sim LN(\mu_i, \sigma^2) \\ 
    \mu_i &= \pmb{X}'\pmb{\beta}\\
    
    \pmb{\beta} &\sim [\pmb{\beta}]
\end{align} \$\$

Here, I've written a general, non-hierarchical case. This approach was
done in detail by Mutshinda et al.~2022, where they modelled the
probability of \(Y_{s,t}\) for different species, \(s\) at times \(t\).
I've rewritten their approach to be consistent with the notation I
described above: \[
\begin{align}
    [y_{s,t}] = \left\{ 
        \begin{array}{l}
            (1 - \pi_{s,t})\\
            \pi_{s,t} \cdot LN(\mu_{s,t}, \sigma^2_s) 
        \end{array} \right.
\end{align}
\]

In this approach, the biomass concentraiton is essentially modeled
separately for each \(s\) species, which has it's own occurance
probability, \(\pi_s\) and expected biomass concentration, \(\mu_s\).
Yet conceivably, it is more appropriate to model similar taxonomic
groups in a hierarchical structure with a group-level mean effect and
species/genus level variation, as I described in approach 0.

\section{Approach 2: Mixed Model}\label{approach-2-mixed-model}

The problem with both approaches described above is that it doesn't
really reflect the true system. Really, plankton occur in the water as
counts - whole individuals. Then the biomass concentration is determined
both by the number of individuals and the size of those individuals:

\[
y_i = n_i \cdot b_i
\]

Observations of the total biomass concentration \(\y_i\) can be modeled
separately as factors contributing to abundance, \(n\) and factors
contributing to the biomass of individuals \(b\).

If we acknowledge that our estimates of \(y_{g,i}\) described previously
are really a derived estimate of counts, \(n\), and a count to biomass
conversion. In many cases, this is a fixed value from some published
study which is then multiplied by the observed counts. However, with the
employment of imaging analyses to process our phytoplankton samples, we
actually are gathering real data on at least the biovolume of each
individual.

So we can really describe
\(y_{g,t} = n_{g,t} \times b_{g,i} \times c_g\) where \(n_{g,t}\) are
the counts of the \(g\) group of phytoplankton at time \(t\). The
biomass of group \(g\) is defined as \(b_{g,i}\) where \(i\) is a subset
of \(t\) where \(n_{g,t}\) \textgreater{} 0 and \(c_g\) is a constant
from literature values.

This approach may offer several advantages. First, by treating our
occurrances as counts, it is possible to include 0s with the
observations, rather than having to develop separate model. This is
particularly useful in the case of variable sampling volume inherent to
imaging studies (or a common net-tow) where the number of individuals
counted may vary based on the sampling effort.

Second, we can also explicitly model the variation around individual
sizes in our second model. By doing this, we can expand the model to
include effects which may influence the size of individual cells. Now we
can ask how both the number and the size of different phytoplankton are
influenced by environmental factors - which have different ecological
implications.

This approach logically matches the system more clearly. Really, the
presence of observing a plankton should follow some count distribution.
Then the size/biomass of individuals are more likely to follow a normal
distribution.

In the first approach we have a lot of derived quantities then try to
model the data but are forced to do weird things like a log-normal
distribution.




\end{document}
